\documentclass[10pt,a4paper]{amsart}
\usepackage[margin=19mm]{geometry}
\usepackage{amsmath,amssymb,mathtools}
\usepackage[scr=rsfs,cal=euler]{mathalfa}
\usepackage{xfrac}
\usepackage[unicode,hidelinks,bookmarks=false]{hyperref}
\newcommand{\ie}{i.e.\ }
\newcommand{\cf}{cf.\ }
\newcommand{\resp}{respectively\ }
\newcommand{\Subsection}{\subsection}
\providecommand{\nobreakdash}{\nobreak\hbox{-}}
\setlength{\emergencystretch}{2em}
\multlinegap0pt
\usepackage{amssymb}
\usepackage{mathtools}
\usepackage{xcolor}
\newtheorem{lemma}{Lemma}
\newcommand{\cP}{\mathcal P}
\newcommand{\dd}{d_{\chi}}
\title{Monochromatic cycle partitions in $3$-mean edge-colourings}
\author{Richard Lang, Guilherme Oliveira Mota}
\date{Source: arXiv:2607.29336v1 (CC BY 4.0)}
\begin{document}
\maketitle

\noindent\emph{Source and licence.} Monochromatic cycle partitions in $3$-mean edge-colourings. Version arXiv:2607.29336v1 (CC BY 4.0), distributed by arXiv under the Creative Commons Attribution 4.0 International licence, http://creativecommons.org/licenses/by/4.0/, as recorded in the arXiv licence metadata for this exact version and shown on the abstract page https://arxiv.org/abs/2607.29336v1. Full licence notice: \texttt{licenses/arxiv-2607.29336-CC-BY-4.0.txt}.\par\smallskip

\noindent\textit{Editorial scope.} Counted unit: the article's lemma lem:bounded-colour-core (source0.tex lines 338--354). The article's complete original proof of it is delivered with it (source0.tex lines 356--406, one \texttt{proof} environment of 51 lines). No mathematics is written, repaired or re-derived: the statement and the proof are the article's own lines, in the article's own notation, and no retained line is substituted or re-flowed.  No further source-local context is needed: every cross-reference of every retained line resolves inside the two delivered spans.  Nothing was omitted from any retained span, so the package carries no omission record.  No retained line contains a citation, so the wrapper carries no bibliography and no bibliography entry.  No retained line contains a figure environment, an image-inclusion command or drawing code, so no image is delivered and the package declares no asset dependency.  Editorial formatting only, in addition: an AMS \texttt{amsart} preamble that restates the article's own declarations and macros used by the retained lines, namely the environments \texttt{lemma} on separate counters, together with the standard packages those lines need.  The retained environments print in this document as Lemma 1 in order of appearance, whereas the article prints the same items with the numbering of its own full document.  The complete native source of the article as distributed in the arXiv e-print for this exact version is bundled unchanged as \texttt{sources/206475/source0.tex}.
\begin{lemma}[Bounded-colour core]\label{lem:bounded-colour-core} Let $H$ be an
  edge-coloured complete graph, and let~$X$ and~$Y'\neq\emptyset$ be disjoint
subsets of $V(H)$. Let $c$ be a colour, $X_0\subseteq X$ and put~$\ell:=\lfloor 5|X|/8\rfloor+1$. If
\begin{align*}
&c\in\cP(u)\text{ and }\dd(u)\le2
  &&\text{for every }u\in X,\quad\text{and}\\
&|X\setminus N_c(v,X)|\ge\ell
  &&\text{for every }v\in Y'\,,
\end{align*}
and
\[
  15|Y'|\le\dd(Y')\le |X|+3|Y'|-12|X_0|\,,
\]
then there is a set $X^*\subseteq X\setminus X_0$ with $|X^*|=|Y'|$ such that
the complete bipartite graph between $X^*$ and $Y'$ uses at most three
colours.
\end{lemma} 

\begin{proof}
Put $x:=|X|$, $y':=|Y'|$, $x_0:=|X_0|$ and $t:=\dd(Y')/y'\geq 15$. For each
$v\in Y'$, we define $t_v:=|\cP(v)\setminus\{c\}|$. Thus, since every vertex of
$Y'$ has a non-$c$ neighbour in $X$, we have $t_v\ge1$, which combined with
$\sum_{v\in Y'}t_v\le\dd(Y')=ty'$ implies that some vertex $v\in Y'$ satisfies
$1\le t_v\le t$.

Let $\sigma:=\ell-x_0$. Deleting $X_0$ removes at most $x_0$ non-$c$ neighbours
of $v$, so $v$ has at least~$\sigma$ non-$c$ neighbours in $X\setminus X_0$.
The key estimate that will allow us to choose the $y'$ vertices forming $X^*$
is
\begin{equation}\label{eq:core-target}
  \frac{2\sigma}{t}>y'\,.
\end{equation}
To see that \eqref{eq:core-target} holds, note first that the bounds on
$\dd(Y')=ty'$ give
$$
x\ge ty'-3y'+12x_0\ge 4ty'/5+12x_0,
$$
where we used that $t\geq 15$ in the last inequality. Thus, since
$\ell>5x/8$, it follows that
$$2\sigma>5x/4-2x_0\ge ty'+13x_0\ge ty',
$$
which proves \eqref{eq:core-target}.

Now let us conclude the proof of the lemma from~\eqref{eq:core-target}. Suppose
first that $t_v=1$. Then the at least $\sigma$ available non-$c$ neighbours
of $v$ belong to the same non-$c$ colour class. By \eqref{eq:core-target} and~$t\ge15$, we have $\sigma>ty'/2>y'$, so this class contains at least $y'$
available vertices. In the case where $t_v\ge2$, write
$d_1\ge\cdots\ge d_{t_v}\ge0$ for
the orders of the corresponding non-$c$ colour classes among these neighbours,
including empty classes if necessary. Since their total order is at least
$\sigma$, we have
\[
  d_1+d_2
  \ge\frac{2}{t_v}\sum_{i=1}^{t_v}d_i
  \ge\frac{2\sigma}{t_v}
  \ge\frac{2\sigma}{t}
  >y'\,,
\]
where the penultimate inequality uses $t_v\le t$. Thus, in either case, the
union of at most two such classes contains at least $y'$ available vertices.

Choose a set $X^*\subseteq X\setminus X_0$ of order $y'$ from these classes.
Let $c_1$ and $c_2$ be their colours; if there is only one such colour, set
$c_2:=c_1$. Thus, for every $u\in X^*$, the edge $uv$ has colour $c_1$ or
$c_2$. This edge is not of colour $c$. Since $c\in\cP(u)$ and
$\dd(u)\le2$, we have
$\cP(u)\subseteq\{c,c_1,c_2\}$. Therefore every edge between $X^*$ and $Y'$
has colour in the set $\{c,c_1,c_2\}$, which has size at most three.
\end{proof}

\end{document}
